% In this file you should put macros to be used only by
% the printed version. Of course they should have a corresponding
% version in macros/web.tex.
% Typically the printed version could have more fancy decorations.
% This should be a very short file.
%
% This file starts with dummy macros that ensure the pdf
% compiler will ignore macros provided by plasTeX that make
% sense only for the web version, such as dependency graph
% macros.


% Dummy macros that make sense only for web version.
\newcommand{\lean}[1]{}
\newcommand{\discussion}[1]{}
\newcommand{\leanok}{}
\newcommand{\mathlibok}{}
\newcommand{\notready}{}
% Print edition renders each TikZ figure from figures/#1 with caption #2.
\newcommand{\blueprintfig}[2]{\begin{figure}[htbp]\centering\input{figures/#1}\caption{#2}\end{figure}}
% Lean names and module paths are set through the url package, which hyperref
% loads: it keeps underscores safe in text mode and, unlike a single \texttt
% box, breaks a long name across lines at the listed characters.
\urlstyle{tt}
\def\UrlBreaks{\do\.\do\_\do\/\do\-}
% A visible Lean declaration reference.
% The two-argument form typesets its short display name alone.
\newcommand{\leandecl}[1]{\nolinkurl{#1}}
\newcommand{\leandeclnamed}[2]{\nolinkurl{#1}}
% Pointer to a whole Lean source file, by slashed module path without extension.
\newcommand{\leanmodule}[1]{\nolinkurl{#1}}
% Make sure that arguments of \uses and \proves are real labels, by using invisible refs:
% latex prints a warning if the label is not defined, but nothing is shown in the pdf file.
% It uses LaTeX3 programming, this is why we use the expl3 package.
\ExplSyntaxOn
\NewDocumentCommand{\uses}{m}
 {\clist_map_inline:nn{#1}{\vphantom{\ref{##1}}}%
  \ignorespaces}
\NewDocumentCommand{\proves}{m}
 {\clist_map_inline:nn{#1}{\vphantom{\ref{##1}}}%
  \ignorespaces}
\ExplSyntaxOff